Make the substitution $y = 1 - \mfrac{1}{8} \cos^2\paren{x}$.
The range of possibe values of $y$ is $\cinterval{1 - \mfrac{1}{8}}{1} = \cinterval{\mfrac{7}{8}}{1}$.
Using that $\sin^2\paren{x} = 1 - \cos^2\paren{x}$ and $\cos^2\paren{x} = 8 - 8 y$,
the substitution yields the equation
\begin{equation*}
y^8 - 8 y + 7 = 0
\end{equation*}
Clearly, $1$ is a solution, and it yields the factorisation
\begin{equation*}
  y^8 - 8 y + 7 = \paren{y - 1}\paren{ - 7 + \sum\limits_{k = 1}^{7} {y^k}}
\end{equation*}
The term $- 7 + \sum\limits_{k = 1}^{7} {y^k}$
is a strictly increasing function of $y$ for $y > 0$.
Thus, in particular, it admits at most one root in the range $\cinterval{\frac{7}{8}}{1}$.
The root is $1$, trivially.
Hence, the equation $y^8 - 8 y + 7 = 0$ admits only the solution $y = 1$ (with multiplicity $2$)
in the range $\cinterval{\frac{7}{8}}{1}$.

The solution $y = 1$ corresponds to $1 - \mfrac{1}{8} \cos^2\paren{x} = 1$,
that is $\cos\paren{x} = 0$.
It is immediate to see also that every $x \in \mathbb{R}$ such that $\cos\paren{x} = 0$
is a solution to the original equation.

In conclusion, the solutions are $x = \mfrac{\pi}{2} + k \pi$ for $k \in \mathbb{Z}$.
